\appendix
\section{Data and Empirical Measurement}
\label{app:empirical-data}

\subsection{Housing stock}

The stock facts use the national public-use household file of the 2023
American Housing Survey, weighted by the housing-unit weight. Tenure has three
occupied categories: owned, rented for cash, and occupied without payment of
cash rent. Shares within a tenure use all units of that tenure as the
denominator; the share of seven-plus-room homes that are rented instead uses
all occupied units of that size. It is 8.2 percent when the denominator
contains owned and cash-rented units and 8.1 percent when it also contains
units occupied without cash rent and units with unreported tenure. Bedrooms and
total rooms are different measures. I report bedrooms for comparison with
other work and rooms because rooms are the PSID and ACS outcome and the unit of
the model. Mean rooms among households with children are computed over
households with at least one child present. A cross-section of occupied homes
does not identify the units offered for rent or the choices families would
make at other prices.

\subsection{PSID first-birth event studies}

\subsubsection*{Sample.}
The PSID data are the family and individual files for 1968--2019. An
observation is a person-year in which the person is a current member of a
family unit and at least 18 years old. The year of first birth is the earliest
birth year of a biological child in the person's relationship-to-child
history, which records up to twenty children. The comparison group consists of
persons whose full history reports no children; persons with incomplete or
unknown histories are dropped. A parent enters only if the first birth
occurs no earlier than the first year in which the person is observed as an
adult in the PSID. All estimates use the PSID longitudinal individual weight
and so include only sample persons with a positive weight.

The household sample keeps women who are the reference person or spouse of a
household containing a single current family unit, and keeps one woman per
household and year, the reference person first. The sample of adults heading a
household at baseline keeps all adults, women and men, who were reference
person or spouse in the window two to three years before the first birth, and
retains all of their person-years. The complementary sample keeps adults who
were in the PSID but not reference person or spouse in that window. Each
sample uses the same comparison group, weights and entry rule.

\subsubsection*{Housing measures.}
All four housing items are taken from the official year-specific PSID family
variables and merged to the interview year in which they were reported. Rooms
are the total rooms in the dwelling. Codes that do not denote a count (9
through 1984; 99 in 1985--1993; 98 and 99 from 1994) are set to missing, and a
reported zero is retained. Ownership equals one if the family owns the
dwelling and zero if it rents; other arrangements are missing. The mobility
item asks whether the family moved since a fixed date: the spring of the
previous survey year in the annual waves, and January 1 of the previous survey
year in recent waves. It therefore records a move within an interval of about
one year in the annual waves and about two years in the biennial waves. The reason for a move is the first-mentioned
reason. ``More space'' is the PSID category for expansion of housing (more
space, more rent, a better place); ``neighborhood'' is the category for
neighborhood-related reasons (a better neighborhood, schooling, proximity to
friends or relatives). Unknown or unreported answers are missing, and code 9
of the reason item (unknown before 2019, homelessness in 2019) is treated as
missing throughout. An observation contributes to an outcome only when that
outcome is reported; unreported outcomes are never coded as zero.

A merged PSID panel used in earlier versions of this analysis stored rooms,
the mobility item and the reason for moving on the row of the interview
before the one at which they were reported, one year earlier through 1997 and
two years earlier afterwards. All 811,486 room values in that panel that
overlap the official variables equal the official value from the following
interview. An earlier estimate of 0.77 rooms three years after the birth
therefore compared rooms four to five years after the birth with rooms in the
year before or the year of the birth; the estimates in the text supersede it.

\subsubsection*{Estimator.}
Event time is the survey year minus the year of first birth, grouped into the
windows $\le-8$, $-7/{-6}$, $-5/{-4}$, $-3/{-2}$, $-1/0$, $+1/{+2}$,
$+3/{+4}$, $+5/{+6}$, $+7/{+8}$, $+9/{+10}$ and $\ge+11$. Two-year windows
are needed because the PSID became biennial after 1997: from the 1999 birth
cohort on, no cohort is observed at both single years $-2$ and $+3$, so a
single-year reference period is observed for only part of the cohorts. The
omitted window is $-3/{-2}$, and each birth cohort must be observed in it;
cohorts without an observation in the reference window are excluded (first
births in 1968--1970 and 1985--1986 in the household sample, none in the
sample fixed at baseline). I estimate cohort-specific event paths relative to
the comparison group and average them with cohort-share weights, following Sun
and Abraham (2021). The regressions include person and survey-year fixed
effects and complete sets of indicators for age and years of education.
Standard errors are clustered by person and account for the estimation of the
cohort shares.

\subsubsection*{Robustness.}
Table~\ref{tab:app-rooms-robust} collects alternative rooms estimates for the
household sample. With the year before the birth in the reference window
($-2/{-1}$), rooms are 0.76 higher two to three years after the birth. Moving
the reference window earlier raises the estimate at three to four years
because it counts the pre-birth rise in this sample; the increase measured from
years $-3/{-2}$ is 1.03, 1.00 and 0.94 across the three reference windows.
Taking rooms from the merged panel, re-dated to the
interview at which they were reported, instead of from the official variables
changes the estimate by 0.02. Using the latest birth cohort instead of the
childless as comparison group leaves it unchanged. Requiring the first birth
to follow the woman's first observation as reference person or spouse, rather
than her first observation as an adult, lowers it by 0.05 and removes 24
percent of the observations.

\begin{table}[H]
\centering
\caption{Rooms around the first birth: alternative samples and reference windows (PSID)}
\label{tab:app-rooms-robust}
\small
\begin{tabularx}{\textwidth}{@{}X l l c r@{}}
\toprule
Sample & Reference & Window & Rooms (SE) & Person-years \\
\midrule
Women reference person or spouse, each year & $-3/{-2}$ & $+3/{+4}$ & 1.03 (0.07) & 63,338 \\
\quad same & $-2/{-1}$ & $+2/{+3}$ & 0.76 (0.06) & 65,053 \\
\quad same & $-5/{-4}$ & $+3/{+4}$ & 1.30 (0.10) & 59,623 \\
\quad same & $-7/{-6}$ & $+3/{+4}$ & 1.43 (0.13) & 56,749 \\
\quad same, first birth after first year as reference person or spouse & $-3/{-2}$ & $+3/{+4}$ & 0.97 (0.07) & 48,225 \\
\addlinespace
Adults reference person or spouse in $-3/{-2}$ & $-3/{-2}$ & $+3/{+4}$ & 1.47 (0.05) & 117,853 \\
Adults not reference person or spouse in $-3/{-2}$ & $-3/{-2}$ & $+3/{+4}$ & $-0.38$ (0.08) & 87,239 \\
All adults & $-3/{-2}$ & $+3/{+4}$ & 0.97 (0.05) & 187,873 \\
\bottomrule
\end{tabularx}
\par\smallskip
\begin{minipage}{0.97\linewidth}
\footnotesize\textit{Notes:} Same estimator, comparison group, controls and
weights as Table~\ref{tab:app-rooms-path}. ``Reference'' is the omitted window
and ``Window'' the reported one, both in years relative to the first birth.
Rows with an earlier reference window measure the change from that window, so
they include the pre-birth rise of the household sample.
\end{minipage}
\end{table}

\subsection{ACS matched pseudo-panel}

The ACS exercise uses the one-year samples for 2005--2019 in all states and
the District of Columbia. The sample consists of women aged 25 to 45. Because
the ACS contains no birth histories, the year of first birth is inferred from
the age of the oldest child linked to the woman in the household roster, and a
woman's housing before that year is represented by childless women surveyed in
the corresponding earlier years who share her demographic matching cell.
Housing outcomes do not enter the matching. Rooms are capped at nine, and codes
that do not denote a count are missing; bedrooms are capped at five; ownership
equals one for owned units and zero for rented ones. The regression includes
event-time indicators from $-5$ to $+10$ with year $-2$ omitted, and state, age
and survey-year fixed effects; it is weighted by the matching weights, with
standard errors clustered by source household because a household can supply
several matched records. The estimation sample has 14.2 million matched
records from 2.8 million source households.

Between one year before and three years after the birth, rooms rise by 0.49
(0.01), bedrooms by 0.34 (0.004), and the ownership rate by 3.3 percentage
points (0.2). Dropping the age fixed effects, or all fixed effects, raises the
room estimate to 0.70 and 0.71 and the ownership estimate to 10.6 and 10.7
points; dropping the state and survey-year effects gives 0.53 rooms and 4.3
points. The pre-birth observations are constructed from other
women, so the design assumes that matched childless women represent the
housing future mothers would have had before the birth. The roster does not
record children living elsewhere or distinguish biological from other
children.

\subsection{Sibling-composition designs}

The ACS sibling designs use women aged 21 to 35 in the 2005--2019 one-year
samples whose oldest linked child, if any, is younger than 18. Children are
linked to mothers through the household roster, which records children living
with the mother, not all children borne, and does not separate biological from
step or adoptive children. The regressions control for complete sets of
indicators for the mother's age at the event, her race, the survey year and
the event age, are weighted by the person weight, and cluster standard errors
by household.

In the first design, the sample is mothers whose oldest linked child is aged 0
to 5 (857,607 mothers). The instrument indicates that the two oldest children
are recorded at the same age, and the treatment is a second resident child.
Ages are recorded in whole years without quarter of birth, so the instrument
is an age-only proxy for a twin birth. The first stage is 0.66 ($F$ statistic
above 77,000). The reduced form is 0.33 rooms (0.02), and the implied effect of
a second child is 0.50 rooms (95 percent interval 0.45 to 0.55), 0.33 bedrooms
(0.31 to 0.36), and 4.7 percentage points in ownership (3.3 to 6.1).

In the second design, the sample is mothers whose two oldest linked children
differ in age and whose second child is aged 0 to 5 (656,986 mothers); it is
not restricted to mothers with exactly two children. The instrument indicates
that the two oldest children have the same sex, and the treatment is a third
resident child. The first stage is 0.033 ($F=715$). Mothers of same-sex pairs
live in 0.038 fewer rooms (0.005) and 0.034 fewer bedrooms (0.003), and their
ownership rate does not differ. Scaled by the first stage, a third child would
reduce the dwelling by 1.15 rooms (95 percent interval 0.83 to 1.47). Children of the same sex can
share a bedroom, so sibling sex composition may change the space a family
needs without changing its size. That possibility would violate the exclusion
restriction for housing outcomes. The data cannot test it formally, and I do
not interpret this estimate as the effect of a third child. The same two designs are too imprecise in the PSID to add
information: same-age births are rare in its smaller sample, and the
sex-composition first stage is weak.

\subsection{Estimates and the second sample}
\label{app:empirical-two-samples}

Tables~\ref{tab:app-rooms-path} and \ref{tab:app-tenure-moves} report the
estimates behind Figures~\ref{fig:emp-rooms-path} and
\ref{fig:emp-tenure-moves}, for the household sample of the main text and for
a second sample that fixes household position before the birth. The second
sample keeps every adult, woman or man, who was a reference person or spouse
two to three years before the birth, and follows that person in every year.
Both partners of a couple can appear in it, so its person-clustered standard
errors are somewhat too small. Figures~\ref{fig:app-rooms-samples} and
\ref{fig:app-tenure-samples} plot the two samples together. Fixing household
position before the birth, the pre-birth path is flat; the household sample's
rise in the four years before the reference window reflects who is in the
sample at that point.

\begin{figure}[p]
\centering
\includegraphics[width=0.9\linewidth]{first_birth_rooms_psid_samples}
\caption{Rooms around the first birth, two samples (PSID)}
\label{fig:app-rooms-samples}
\end{figure}

\begin{figure}[p]
\centering
\includegraphics[width=0.95\linewidth]{first_birth_tenure_moves_psid_samples}
\caption{Ownership and moving around the first birth, two samples (PSID)}
\label{fig:app-tenure-samples}
\end{figure}

\begin{table}[p]
\centering
\caption{Rooms around the first birth (PSID)}
\label{tab:app-rooms-path}
\small
\begin{tabular}{@{}l c c@{}}
\toprule
 & Women who are reference & Adults who were reference \\
 & person or spouse & person or spouse two to \\
 & in each year & three years before the birth \\
\midrule
\multicolumn{3}{@{}l}{\textit{Years relative to first birth: change in rooms relative to years $-3$ and $-2$}} \\
$-8$ or earlier & $-0.50$ (0.12) & $\phantom{-}0.15$ (0.08) \\
$-7$ and $-6$ & $-0.50$ (0.10) & $\phantom{-}0.04$ (0.06) \\
$-5$ and $-4$ & $-0.34$ (0.06) & $-0.03$ (0.04) \\
$-3$ and $-2$ & 0 & 0 \\
$-1$ and $0$ & $\phantom{-}0.41$ (0.05) & $\phantom{-}0.61$ (0.04) \\
$+1$ and $+2$ & $\phantom{-}0.82$ (0.06) & $\phantom{-}1.19$ (0.04) \\
$+3$ and $+4$ & $\phantom{-}1.03$ (0.07) & $\phantom{-}1.47$ (0.05) \\
$+5$ and $+6$ & $\phantom{-}1.17$ (0.08) & $\phantom{-}1.62$ (0.06) \\
$+7$ and $+8$ & $\phantom{-}1.20$ (0.09) & $\phantom{-}1.64$ (0.07) \\
$+9$ and $+10$ & $\phantom{-}1.25$ (0.10) & $\phantom{-}1.83$ (0.08) \\
$+11$ or later & $\phantom{-}1.16$ (0.13) & $\phantom{-}1.85$ (0.09) \\
\midrule
Mean rooms, years $-3$ and $-2$ & 4.70 & 4.69 \\
Person-years & 63,338 & 117,853 \\
Persons: first-time parents / childless & 3,655 / 1,654 & 3,302 / 6,008 \\
\bottomrule
\end{tabular}
\par\smallskip
\begin{minipage}{0.97\linewidth}
\footnotesize\textit{Notes:} Interaction-weighted event-study estimates
(Sun and Abraham, 2021) in two-year windows of event time, with standard errors
clustered by person in parentheses. The comparison group consists of adults
with no children in their full relationship history. All regressions include
person and survey-year fixed effects and indicators for age and years of
education, and use PSID longitudinal weights. Rooms are the total rooms in the
dwelling reported at the interview. The first column excludes first births in
1968--1970 and 1985--1986, which have no observation in the reference window.
\end{minipage}
\end{table}

\begin{table}[p]
\centering
\caption{Ownership and moving around the first birth (PSID)}
\label{tab:app-tenure-moves}
\footnotesize
\setlength{\tabcolsep}{4pt}
\begin{tabular}{@{}l c c c c c r@{}}
\toprule
 & Mean, & \multicolumn{4}{c}{Change relative to years $-3$ and $-2$} & Person- \\
\cmidrule(lr){3-6}
 & $-3/{-2}$ & $-1/0$ & $+1/{+2}$ & $+3/{+4}$ & $+5/{+6}$ & years \\
\midrule
\multicolumn{7}{@{}l}{\textit{A. Women who are reference person or spouse in each year}} \\
Owns the dwelling & 0.41 & $\phantom{-}0.11$ (0.01) & $\phantom{-}0.16$ (0.02) & $\phantom{-}0.16$ (0.02) & $\phantom{-}0.14$ (0.02) & 61,229 \\
Moved since last interview & 0.58 & $-0.04$ (0.02) & $-0.12$ (0.02) & $-0.16$ (0.02) & $-0.16$ (0.02) & 63,817 \\
Moved for more space & 0.070 & $\phantom{-}0.014$ (0.009) & $\phantom{-}0.034$ (0.009) & $\phantom{-}0.020$ (0.010) & $\phantom{-}0.007$ (0.010) & 62,813 \\
Moved for neighborhood & 0.032 & $-0.004$ (0.006) & $-0.005$ (0.006) & $-0.010$ (0.005) & $\phantom{-}0.002$ (0.006) & 62,813 \\
\multicolumn{7}{@{}l}{\quad\textit{Share of moves, households that moved:}} \\
\quad Share for more space & 0.122 & $\phantom{-}0.054$ (0.017) & $\phantom{-}0.140$ (0.021) & $\phantom{-}0.133$ (0.023) & $\phantom{-}0.108$ (0.026) & 19,456 \\
\quad Share for neighborhood & 0.054 & $-0.017$ (0.012) & $-0.011$ (0.014) & $-0.017$ (0.014) & $\phantom{-}0.022$ (0.018) & 19,456 \\
\addlinespace
\multicolumn{7}{@{}l}{\textit{B. Adults who were reference person or spouse two to three years before the birth}} \\
Owns the dwelling & 0.43 & $\phantom{-}0.16$ (0.01) & $\phantom{-}0.25$ (0.01) & $\phantom{-}0.27$ (0.01) & $\phantom{-}0.26$ (0.01) & 111,272 \\
Moved since last interview & 0.54 & $-0.08$ (0.01) & $-0.17$ (0.01) & $-0.21$ (0.01) & $-0.21$ (0.01) & 118,875 \\
Moved for more space & 0.064 & $\phantom{-}0.015$ (0.006) & $\phantom{-}0.021$ (0.006) & $\phantom{-}0.010$ (0.006) & $\phantom{-}0.009$ (0.006) & 116,697 \\
Moved for neighborhood & 0.030 & $-0.009$ (0.004) & $-0.008$ (0.004) & $-0.009$ (0.004) & $-0.003$ (0.004) & 116,697 \\
\multicolumn{7}{@{}l}{\quad\textit{Share of moves, households that moved:}} \\
\quad Share for more space & 0.123 & $\phantom{-}0.050$ (0.013) & $\phantom{-}0.114$ (0.017) & $\phantom{-}0.142$ (0.020) & $\phantom{-}0.147$ (0.022) & 27,462 \\
\quad Share for neighborhood & 0.058 & $-0.022$ (0.009) & $\phantom{-}0.005$ (0.011) & $-0.002$ (0.013) & $\phantom{-}0.012$ (0.015) & 27,462 \\
\bottomrule
\end{tabular}
\par\smallskip
\begin{minipage}{0.97\linewidth}
\footnotesize\textit{Notes:} Same estimator, comparison group, controls,
weights and reference window as Table~\ref{tab:app-rooms-path}. Columns are
windows of years relative to the first birth. Entries are changes relative to
years $-3$ and $-2$, with person-clustered standard errors
in parentheses; the first column is the weighted mean in the reference window.
Each outcome uses its own sample of non-missing observations. ``Moved for more
space'' equals one if the household moved since the previous interview and
gave expansion of housing (more space, a better place) as the first reason,
and zero otherwise; ``moved for neighborhood'' is defined in the same way. The
rows for households that moved restrict the sample to household-years with a
move and therefore condition on an outcome of the birth.
\end{minipage}
\end{table}
